\documentclass[11pt]{article}
\usepackage[margin=1in]{geometry}
\usepackage{booktabs}
\usepackage{amsmath}
\usepackage{fancyhdr}
\usepackage{xcolor}
\usepackage{hyperref}

\pagestyle{fancy}
\fancyhf{}
\lhead{Barrier options under Heston SV: Python replication}
\rhead{\thepage}

\newcommand{\good}{\textcolor{green!45!black}{good}}
\newcommand{\fair}{\textcolor{orange!70!black}{fair}}

\title{Pricing Barrier Options under Heston Stochastic Volatility:\\
A Python Implementation of the Method of Lines and 2D COS Method,\\
Cross-Validated against Monte Carlo and Chiarella, Kang \& Meyer (2012)\\[0.3em]
\large Now Extended to American Exercise via Bermudan-COS}
\author{Replication report}
\date{\today}

\begin{document}
\maketitle

\begin{abstract}
This note reports on a Python implementation of independent pricing engines
for up-and-out barrier call options under Heston stochastic volatility
(optionally with Merton/Bates log-normal jumps): a \emph{method-of-lines}
(MOL) finite-difference PDE solver ported from the Fortran reference code, a
\emph{2D COS} (Fourier-cosine) expansion pricer ported from the MATLAB
reference code, an independent \emph{Monte Carlo} (MC) simulator, and a new
\emph{Bermudan-COS American} extension of the 2D COS pricer built for this
report. All are cross-validated against each other and against the published
numerical examples of Chiarella, Kang \& Meyer (2012, \emph{Comput. Math.
Appl.} 64, 2034--2048) -- the paper this codebase's method of lines
implements. We report replication accuracy, discuss a grid-sensitivity issue
found in the MOL discrete-monitoring code path, and present the new
COS-American engine: its numerical design, two genuine numerical
instabilities found and fixed during development, and its accuracy against
MOL for discretely- and continuously-monitored American barrier options,
with and without jumps.
\end{abstract}

% Table numbering below is deliberately offset by 1 (first local table is
% "Table 2") to match Chiarella, Kang & Meyer (2012)'s own table numbers,
% since their Table 1 (the parameter table) is referenced in prose here but
% not reproduced as a table of its own.
\setcounter{table}{1}

\section{The three (now four) engines}

\begin{tabular}{@{}p{0.16\textwidth}p{0.8\textwidth}@{}}
\texttt{mol.py} & Finite-difference method of lines: the Heston PDE is
discretized along lines of constant variance, the resulting 2nd-order ODE
system in the spot variable is solved via a Riccati transform
(forward/backward sweep), and the variance lines are coupled by a
Gauss--Seidel sweep, all inside an implicit-Euler-then-BDF2 time stepping
loop. Handles continuous or discrete monitoring, European or American
exercise, and knock-out or (via in-out parity) knock-in payoffs. Price, delta
and gamma all come directly out of the ODE system, with no finite
differencing. The per-line recursion is Numba-JIT compiled. \\[0.5em]
\texttt{cos2d.py} & 2D Fourier-cosine expansion: variance is discretized on
$M{=}200$ Clenshaw--Curtis quadrature nodes; for each node, a standard
discretely-monitored 1D COS recursion (Fang \& Oosterlee 2009) runs in the
log-price dimension \emph{as if} variance were frozen at that node for the
remaining life of the option, and only the very last step blends the $M$
frozen-variance branches together using the true one-step transition density
of variance from today's $v_0$. European, knock-out only, discrete monitoring
only. \\[0.5em]
\texttt{cos2d\_american.py} & \textbf{New.} Bermudan-COS extension of the
same frozen-variance-branch architecture to American exercise, both discrete-
and continuous-barrier monitoring. See Section~4. \\[0.5em]
\texttt{mc.py} & Monte Carlo simulation of the same Heston(+jump) dynamics
(full-truncation Euler scheme for the variance), with a Broadie--Glasserman--Kou
continuity correction applied only where the barrier is genuinely continuously
monitored. Delta/gamma via common-random-number central bumps. European
only. \\
\end{tabular}

All four share a common parameter/spec layer (\texttt{params.py}) so that the
same option and Heston(+jump) specification can be priced by whichever subset
of engines supports it.

\section{Replication of Chiarella, Kang \& Meyer (2012)}

Parameters throughout (Table 1 of the paper): $r{=}0.03$, $q{=}0.05$,
$T{=}0.5$, $K{=}100$, $H{=}130$, $\kappa_v{=}2.00$, $\theta{=}0.10$,
$\sigma{=}0.10$, $\rho{=}-0.50$, $v_0{=}0.10$, no jumps. Up-and-out call,
spot $S_0 \in \{80,90,100,110,120\}$.

\subsection*{A discrepancy in the paper's own description, resolved by triangulation}

The paper's Section~4 states discretely-monitored examples use ``weekly
monitoring frequency''. Taken literally ($\approx$26 evenly-spaced dates over
$T{=}0.5$y), \emph{none} of the engines reproduce Tables~4--5 -- MOL, COS and
MC all independently priced well below the published values, and a
Broadie--Glasserman--Kou continuity-correction check confirmed that scale of
price is what $\sim$26 monitoring dates should actually produce. What
\emph{does} reproduce Tables~4--5 almost exactly is monitoring at just the
schedule hardcoded in the companion Fortran solvers already in this
repository: 2 dates ($T/2, T$) for the European case (\texttt{Eur\_disc.f90})
and 3 ($T/3, 2T/3, T$) for the American case (\texttt{Ame\_disc.f90}). We use
that schedule below.

\begin{table}[h]
\centering
\caption{Continuous monitoring, European (no early exercise)}
\begin{tabular}{@{}rrrr@{}}
\toprule
$S_0$ & MOL (this work) & MOL (paper, finest grid) & abs.\ diff \\
\midrule
80  & 0.9038 & 0.9046 & 0.0008 \\
90  & 1.8747 & 1.8806 & 0.0059 \\
100 & 2.5854 & 2.5999 & 0.0145 \\
110 & 2.4709 & 2.4858 & 0.0149 \\
120 & 1.4740 & 1.4812 & 0.0072 \\
\bottomrule
\end{tabular}
\end{table}

Grid: $n_S{=}150$, $M{=}150$, 150--300 time steps. \good{} match throughout.

\begin{table}[h]
\centering
\caption{Continuous monitoring, American (early exercise) -- \textbf{now with COS}}
\begin{tabular}{@{}rrrrr@{}}
\toprule
$S_0$ & COS (py, new) & MOL (this work) & MOL (paper, finest grid) & COS vs.\ paper \\
\midrule
80  & 1.5095 & 1.4029 & 1.4012 & $+7.73\%$ \\
90  & 4.0346 & 3.9383 & 3.9364 & $+2.50\%$ \\
100 & 8.3199 & 8.3051 & 8.3003 & $+0.24\%$ \\
110 & 14.3002 & 14.4206 & 14.4033 & $-0.72\%$ \\
120 & 21.5820 & 21.8732 & 21.8219 & $-1.10\%$ \\
\bottomrule
\end{tabular}
\end{table}

MOL: \good{} match ($<0.3\%$ relative throughout), including the free-boundary /
early-exercise machinery. \textbf{COS (new):} excellent at-the-money and
in-the-money ($<1.1\%$ from the paper, $S_0 \ge 100$); a real, consistently
$\sim$constant-absolute ($\approx\$0.10$) bias remains for deep out-of-the-money
spots, discussed in Section~4.

\begin{table}[h]
\centering
\caption{Discrete monitoring (2 dates: $T/2, T$), European}
\begin{tabular}{@{}rrrrrr@{}}
\toprule
$S_0$ & COS (py) & COS (paper) & MOL (py) & MOL (paper) & MC (py) \\
\midrule
80  & 1.0805 & 1.0809 & 1.0659 & 1.0807 & 1.0557 \\
90  & 2.4848 & 2.4871 & 2.4613 & 2.5289 & 2.4401 \\
100 & 4.0405 & 4.0454 & 3.9474 & 4.1116 & 3.9728 \\
110 & 4.9710 & 4.9779 & 4.7517 & 5.0235 & 4.7723 \\
120 & 4.8578 & 4.8646 & 4.5329 & 4.8706 & 4.5563 \\
\bottomrule
\end{tabular}
\end{table}

\textbf{COS}: \good{}, max abs.\ diff 0.007 ($<0.2\%$) -- essentially exact.
\textbf{MOL / MC}: \fair{}, a consistent 2--4\% low bias that does not shrink
cleanly with grid/path refinement (investigated in Section~3).

\begin{table}[h]
\centering
\caption{Discrete monitoring (3 dates: $T/3, 2T/3, T$), American -- \textbf{now with COS}}
\begin{tabular}{@{}rrrrr@{}}
\toprule
$S_0$ & COS (py, new) & MOL (this work) & MOL (paper) & COS vs.\ paper \\
\midrule
80  & 1.5123 & 1.4120 & 1.4008 & $+7.97\%$ \\
90  & 4.0439 & 3.9512 & 3.9339 & $+2.80\%$ \\
100 & 8.3648 & 8.3256 & 8.3010 & $+0.77\%$ \\
110 & 14.4605 & 14.4826 & 14.4446 & $+0.11\%$ \\
120 & 22.0256 & 22.0926 & 22.0389 & $-0.06\%$ \\
\bottomrule
\end{tabular}
\end{table}

MOL: \good{} match ($<0.4\%$ relative), despite using the same
discrete-monitoring domain-switch code as Table~4's European case, which is
the puzzle addressed next. \textbf{COS (new)}: the same pattern as the
continuous case above -- excellent at $S_0 \ge 100$ ($<0.8\%$), a real
$\sim$8\% relative (but modest, $\approx\$0.11$, absolute) gap deep
out-of-the-money. The paper itself has no COS column here (it only supports
European exercise in the original MATLAB code); this is a new comparison, not
a discrepancy against a published number.

\section{Investigating the MOL European discrete-monitoring bias}

Table~4's MOL/MC $\sim$2--4\% low bias, and its failure to shrink
monotonically under grid refinement ($n_S{=}75/150/250$ gave 4.02/3.95/4.03 at
$S_0{=}100$), was investigated directly rather than left unexplained.

\begin{enumerate}
\item \textbf{Ruled out -- SOR/Gauss--Seidel under-convergence.} Tightening
the sweep tolerance from $10^{-8}$ to $10^{-10}$ and the iteration cap from 60
to 150 changed the price by less than $10^{-4}$: the per-step nonlinear solve
is fully converged; this is not the source.
\item \textbf{Ruled out -- a monitoring-date/tau convention mismatch.} MOL,
COS and MC each convert the shared \texttt{OptionSpec.monitor\_dates} field
via the same $T - \texttt{monitor\_dates}$ expression, but MOL's own internal
time axis runs the \emph{opposite} direction ($\tau{=}0$ at maturity, $\tau{=}T$
``now'') from the axis COS/MC use. Testing the corrected input directly
reproduced the \emph{identical} prices to four decimal places. This is not
the source (though the convention split is real and worth documenting to
avoid future confusion).
\item \textbf{Leading hypothesis -- domain-switch remeshing sensitivity,
masked by intrinsic value in the American case.} At each discrete monitoring
date, MOL remeshes the price/delta surface from the wide domain $(0,S_{\max})$
onto the narrow domain $(0,H)$ and back via linear interpolation, against a
sinh-stretched grid whose exact point placement near $S{=}H$ shifts
non-smoothly as $n_S$ changes. The resulting interpolation error appears
comparable in \emph{absolute} size for European and American contracts, but
the European price here ($\approx 4$) has no intrinsic-value floor to absorb
it, whereas the American price ($\approx 8.3$) is dominated by the robust,
purely-algebraic exercise region -- so the same absolute noise reads as
2--4\% for European and $<0.4\%$ for American. This is consistent with all
observations above but was not converted into a code fix in this pass; a
natural next step is a higher-order (e.g.\ cubic) remesh, or local grid
refinement pinned exactly at $S{=}H$ independent of $n_S$.
\end{enumerate}

\section{The new COS-American engine (\texttt{cos2d\_american.py})}

\texttt{cos2d.py}'s discretely-monitored recursion never leaves Fourier space
between monitoring dates -- the barrier is enforced by restricting an
FFT-convolution kernel to the sub-interval $(a,h)$, a trick that exists
specifically to avoid ever evaluating the option value pointwise. American
exercise has no such trick available in general: $\max(\text{continuation},
\text{intrinsic})$ is not a linear operation representable as a convolution.
Following Fang \& Oosterlee (2009), the new module leaves Fourier space at
every exercise opportunity -- evaluate the continuation value pointwise on a
grid, take the max with the intrinsic payoff, and numerically re-expand the
result back onto Fourier-cosine coefficients (a forward cosine transform,
implemented as an explicit real cosine/sine basis matrix rather than an FFT,
for a self-verifiable round trip) before continuing backward. Because American
exercise is not naturally discrete the way a monitoring date is, the engine
always runs a dense Bermudan grid of \texttt{n\_exercise} exercise dates and
Richardson-extrapolates across increasing \texttt{n\_exercise} to approximate
continuous exercise; ``discrete'' vs.\ ``continuous'' \emph{barrier}
monitoring is then just a choice of which of those dates also enforce the
knock-out truncation.

\subsection*{Two numerical instabilities found and fixed during development}

\begin{enumerate}
\item \textbf{Chaining the true Heston transition is not a semigroup.}
\texttt{cos2d.py}'s frozen-variance branches apply the true one-period
Heston transition characteristic function once or twice per branch
($mt{-}1 \le 2$ in the validated tables above). Chaining that same formula
tens of times (needed for a dense Bermudan grid), always re-anchoring to the
same starting variance at every fine step, is not a semigroup:
$\phi_{\text{Heston}}(u,v,dt)^n \ne \phi_{\text{Heston}}(u,v,n\,dt)$ (verified
numerically -- the gap saturates but is non-negligible), and iterating this
through the evaluate/reproject round trip tens of times compounded into a
$>15\%$ price bias by 50 steps, growing with the step count. \textbf{Fix:}
each branch's fine recursion uses a constant-variance Black--Scholes(+jump)
characteristic function instead -- an exact L\'evy semigroup, verified
floating-point-exact under chaining, and arguably the more literal reading
of ``variance frozen for the whole remaining life'' anyway.
\item \textbf{The final cross-branch blend's quadrature breaks down at fine
$dt$.} \texttt{cos2d.py}'s one-shot blend weights branches by the true
one-step variance-transition density, evaluated by quadrature over its
$M$-node Clenshaw--Curtis grid -- valid because its own $dt{=}T/mt$ is large.
Reusing that blend at this module's tiny fine-step $dt$ was tried first: the
transition density's width shrinks below the quadrature grid's local node
spacing, its normalization drifts from $\approx 1$ to $>2$ as $dt$ shrinks,
and the resulting price swung non-monotonically by $\sim$10\% from changing
$M$ alone (75 vs.\ 300), independent of $N$ or \texttt{n\_exercise}.
\textbf{Fix:} since each branch's reconstructed value is already smooth in
variance (verified: no node-to-node noise across dense $M$ sweeps), and since
variance cannot have moved from $v_0$ over an infinitesimal step anyway, the
final read-out is plain interpolation of the (already exercise/barrier-corrected)
per-branch values at the real $v_0$ -- both simpler and the correct $dt\to0$
limit. Verified to remove the $M$-sensitivity entirely (identical to four
significant figures across $M{=}75$--$300$).
\item \textbf{Continuous monitoring needs $N$ to scale with
\texttt{n\_exercise}.} Discrete barrier schedules truncate the value to zero
above $h$ at only 2--3 fixed dates, so the number of times a fresh
discontinuity is reprojected through the $N$-term cosine basis stays small
even as \texttt{n\_exercise} grows -- $N{=}256$ converges cleanly there into
the hundreds of exercise dates. Continuous monitoring truncates at
\emph{every} exercise node, so that cost scales with \texttt{n\_exercise}
itself: at $N{=}256$ the $S_0{=}120$ continuous case converges smoothly only
up to $\texttt{n\_exercise}\approx320$ before visibly diverging (21.83 at
320 $\to$ 23.18 at 640 $\to$ 23.92 at 1280, versus a converged value near
21.8). Doubling $N$ to 512 pushed the divergence onset comfortably past the
\texttt{n\_exercise} range needed for good accuracy. This is documented as a
usage requirement in the module (not solved with an automatic guard):
continuously-monitored runs need a larger $N$ than discretely-monitored ones
at the same \texttt{n\_exercise}.
\end{enumerate}

None of these are subtle -- each was caught by a direct numerical experiment
(chain the transition analytically and compare; sweep $M$ alone holding
everything else fixed; sweep \texttt{n\_exercise} alone and watch for
monotonic convergence) before being accepted as fixed, following the same
``trust but verify against a structurally unrelated method'' discipline
Section~3 above uses for MOL.

\section{American exercise with jumps}

Merton/Bates log-normal jump parameters $\lambda{=}0.5$, $\gamma{=}-0.05$,
$\delta{=}0.1$ (same illustrative values as the package's own quick-start
example), otherwise the Table 3/5 parameters above. \texttt{cos2d\_american.py}
inherits \texttt{cos2d.py}'s jump characteristic-function convention (the
log-jump-size mean is shifted \emph{inside} the CF itself, no separate drift
compensator), which is \emph{not} interchangeable with \texttt{charfunc.py}'s
compensated convention that \texttt{mol.py} uses -- see \texttt{cos2d.py}'s
docstring for why (a calibration-locked choice, not something to "fix"). So
the comparison below validates that the American-exercise/barrier machinery
behaves correctly with jumps active -- sensible, stable numbers, same
accuracy pattern as the no-jump case -- not exact cross-engine agreement,
which the convention difference alone would preclude regardless of any other
error.

\begin{table}[h]
\centering
\caption{American barrier with jumps ($\lambda{=}0.5,\gamma{=}-0.05,\delta{=}0.1$)}
\begin{tabular}{@{}rrrr@{}}
\toprule
$S_0$ & COS+jump (py) & MOL+jump (py) & diff \\
\midrule
\multicolumn{4}{@{}l}{\emph{Discrete monitoring (3 dates: $T/3,2T/3,T$)}} \\
80  & 1.4686 & 1.3840 & $+6.11\%$ \\
100 & 8.1009 & 8.1049 & $-0.05\%$ \\
120 & 21.5928 & 21.7190 & $-0.58\%$ \\
\multicolumn{4}{@{}l}{\emph{Continuous monitoring}} \\
80  & 1.4637 & 1.3780 & $+6.22\%$ \\
100 & 8.0336 & 8.0841 & $-0.62\%$ \\
120 & 21.2296 & 21.5864 & $-1.65\%$ \\
\bottomrule
\end{tabular}
\end{table}

The pattern is the same one seen throughout Section~4: excellent agreement
at- and in-the-money, a larger (but stable, non-diverging) gap deep
out-of-the-money, on top of the expected small residual from the jump-CF
convention difference. Nothing here suggests the jump term interacts badly
with the exercise/barrier machinery -- the deep-OTM gap size and direction
are consistent with the no-jump case, not amplified by adding jumps.

\section{Summary}

The COS method, previously restricted to European discretely-monitored
barrier options in this codebase, now also prices American exercise
(discrete or continuous barrier, with or without jumps) via a Bermudan-COS
recursion, at a level of accuracy (excellent ATM/ITM, a modest deep-OTM bias)
and speed (still faster than MOL, though not by the $\sim$100$\times$ margin
of the pure-European case, since American exercise requires leaving Fourier
space at every one of many exercise dates rather than a handful of monitoring
dates) that makes it a genuine second, structurally-independent method for
American barrier options under Heston(+jumps) -- alongside, not a
replacement for, \texttt{mol.py}, which remains the more accurate choice
deep out-of-the-money until that residual bias is root-caused.

\end{document}
